\documentclass[10pt,a4paper]{amsart}
\usepackage{amsmath,amssymb,amsthm,graphicx}
\usepackage{verbatim}
\newenvironment{reference}{\comment}{\endcomment}
%\newenvironment{reference}{}{}
\newenvironment{slogan}{\comment}{\endcomment}
\newenvironment{history}{\comment}{\endcomment}

% For commutative diagrams we use Xy-pic
\usepackage[all]{xy}

% We use 2cell for 2-commutative diagrams.
\xyoption{2cell}
% Theorem environments.
%
\theoremstyle{plain}
\newtheorem{theorem}[subsection]{Theorem}
\newtheorem{proposition}[subsection]{Proposition}
\newtheorem{lemma}[subsection]{Lemma}

\theoremstyle{definition}
\newtheorem{definition}[subsection]{Definition}
\newtheorem{example}[subsection]{Example}
\newtheorem{exercise}[subsection]{Exercise}
\newtheorem{situation}[subsection]{Situation}

\theoremstyle{remark}
\newtheorem{remark}[subsection]{Remark}
\newtheorem{remarks}[subsection]{Remarks}

\numberwithin{equation}{subsection}

% Macros
%
\def\lim{\mathop{\mathrm{lim}}\nolimits}
\def\colim{\mathop{\mathrm{colim}}\nolimits}
\def\Spec{\mathop{\mathrm{Spec}}}
\def\Hom{\mathop{\mathrm{Hom}}\nolimits}
\def\Ext{\mathop{\mathrm{Ext}}\nolimits}
\def\SheafHom{\mathop{\mathcal{H}\!\mathit{om}}\nolimits}
\def\SheafExt{\mathop{\mathcal{E}\!\mathit{xt}}\nolimits}
\def\Sch{\mathit{Sch}}
\def\Mor{\mathop{\mathrm{Mor}}\nolimits}
\def\Ob{\mathop{\mathrm{Ob}}\nolimits}
\def\Sh{\mathop{\mathit{Sh}}\nolimits}
\def\NL{\mathop{N\!L}\nolimits}
\def\CH{\mathop{\mathrm{CH}}\nolimits}
\def\proetale{{pro\text{-}\acute{e}tale}}
\def\etale{{\acute{e}tale}}
\def\QCoh{\mathit{QCoh}}
\def\Ker{\mathop{\mathrm{Ker}}}
\def\Im{\mathop{\mathrm{Im}}}
\def\Coker{\mathop{\mathrm{Coker}}}
\def\Coim{\mathop{\mathrm{Coim}}}

% Boxtimes
%
\DeclareMathSymbol{\boxtimes}{\mathbin}{AMSa}{"02}

%
% Macros for moduli stacks/spaces
%
\def\QCohstack{\mathcal{QC}\!\mathit{oh}}
\def\Cohstack{\mathcal{C}\!\mathit{oh}}
\def\Spacesstack{\mathcal{S}\!\mathit{paces}}
\def\Quotfunctor{\mathrm{Quot}}
\def\Hilbfunctor{\mathrm{Hilb}}
\def\Curvesstack{\mathcal{C}\!\mathit{urves}}
\def\Polarizedstack{\mathcal{P}\!\mathit{olarized}}
\def\Complexesstack{\mathcal{C}\!\mathit{omplexes}}
% \Pic is the operator that assigns to X its picard group, usage \Pic(X)
% \Picardstack_{X/B} denotes the Picard stack of X over B
% \Picardfunctor_{X/B} denotes the Picard functor of X over B
\def\Pic{\mathop{\mathrm{Pic}}\nolimits}
\def\Picardstack{\mathcal{P}\!\mathit{ic}}
\def\Picardfunctor{\mathrm{Pic}}
\def\Deformationcategory{\mathcal{D}\!\mathit{ef}}
\usepackage{cleveref}
\crefname{originalreference}{original reference}{original references}
\Crefname{originalreference}{Original reference}{Original references}
\makeatletter
\newlabel{homology-definition-serre-subcategory}{{02MO}{source}{}{sourcepointer-1}{}}
\newlabel{homology-definition-serre-subcategory@cref}{{[originalreference][0][]02MO}{[1][0][]source}{}{}{}}
\makeatother
\crefname{originalreference}{original reference}{original references}
\Crefname{originalreference}{Original reference}{Original references}
\makeatletter
\newlabel{stacks-sheaves-lemma-base-change-higher-direct-images}{{075H}{source}{}{sourcepointer-2}{}}
\newlabel{stacks-sheaves-lemma-base-change-higher-direct-images@cref}{{[originalreference][0][]075H}{[1][0][]source}{}{}{}}
\makeatother
\crefname{originalreference}{original reference}{original references}
\Crefname{originalreference}{Original reference}{Original references}
\makeatletter
\newlabel{stacks-sheaves-proposition-smooth-covering-compute-direct-image}{{072G}{source}{}{sourcepointer-3}{}}
\newlabel{stacks-sheaves-proposition-smooth-covering-compute-direct-image@cref}{{[originalreference][0][]072G}{[1][0][]source}{}{}{}}
\makeatother
\crefname{originalreference}{original reference}{original references}
\Crefname{originalreference}{Original reference}{Original references}
\makeatletter
\newlabel{homology-lemma-first-quadrant-ss}{{0132}{source}{}{sourcepointer-4}{}}
\newlabel{homology-lemma-first-quadrant-ss@cref}{{[originalreference][0][]0132}{[1][0][]source}{}{}{}}
\makeatother
\crefname{originalreference}{original reference}{original references}
\Crefname{originalreference}{Original reference}{Original references}
\makeatletter
\newlabel{stacks-morphisms-lemma-quasi-compact-quasi-separated-permanence}{{075S}{source}{}{sourcepointer-5}{}}
\newlabel{stacks-morphisms-lemma-quasi-compact-quasi-separated-permanence@cref}{{[originalreference][0][]075S}{[1][0][]source}{}{}{}}
\makeatother
\title{Higher direct images preserve a smooth-local weak Serre class for qcqs morphisms}
\author{The Stacks Project Authors}
\date{stacks-project:stacks-cohomology.tex:local-snapshot-20261012; GNU FDL 1.2 or later}
\begin{document}
\maketitle
\noindent\textit{Editorial note.} All delivered mathematical text is the author's own, copied verbatim from the pinned Stacks Project chapter file stacks-cohomology.tex; no statement or proof step was rewritten, reordered or removed.\par\medskip
\noindent\textit{Editorial note.} The counted span cites two Sheaves on Stacks results, one Homology lemma and one Morphisms of Stacks lemma, recorded in assembly.external\_references with canonical tags as navigation pointers.\par\medskip
\noindent\textit{Allowed standard prerequisites (editorial).}\begin{itemize}
\item Base change of higher direct images. Stacks Project, Sheaves on Stacks, Lemma stacks-sheaves-lemma-base-change-higher-direct-images. Supports the first chart replacement.
\item Smooth covering spectral sequence. Stacks Project, Sheaves on Stacks, Proposition stacks-sheaves-proposition-smooth-covering-compute-direct-image. Computes \$R f\_*\$ from the \$\textbackslash{}check\{\textbackslash{}mathrm\{C\}\}\$ech terms of the smooth cover.
\item First quadrant spectral sequences. Stacks Project, Homology, Lemma homology-lemma-first-quadrant-ss. Supplies the convergence property used.
\item qcqs permanence for algebraic stacks. Stacks Project, Morphisms of Stacks, Lemma stacks-morphisms-lemma-quasi-compact-quasi-separated-permanence. Keeps the \$\textbackslash{}mathcal\{V\}\_p\$ qcqs through the iterations.
\end{itemize}\medskip
\ifdefined\hypertarget\hypertarget{sourcepointer-1}{}\fi
\noindent\textit{Source reference.} Original reference 02MO: sources/stacks-chow/tags.txt, source lines 3451--3451. Cross-chapter standard prerequisite cited by the author inside the counted span; it is used as a pointer and its own proof is not part of this item.\par\medskip
\ifdefined\hypertarget\hypertarget{sourcepointer-2}{}\fi
\noindent\textit{Source reference.} Original reference 075H: sources/stacks-chow/tags.txt, source lines 9131--9131. Cross-chapter standard prerequisite cited by the author inside the counted span; it is used as a pointer and its own proof is not part of this item.\par\medskip
\ifdefined\hypertarget\hypertarget{sourcepointer-3}{}\fi
\noindent\textit{Source reference.} Original reference 072G: sources/stacks-chow/tags.txt, source lines 9025--9025. Cross-chapter standard prerequisite cited by the author inside the counted span; it is used as a pointer and its own proof is not part of this item.\par\medskip
\ifdefined\hypertarget\hypertarget{sourcepointer-4}{}\fi
\noindent\textit{Source reference.} Original reference 0132: sources/stacks-chow/tags.txt, source lines 1420--1420. Cross-chapter standard prerequisite cited by the author inside the counted span; it is used as a pointer and its own proof is not part of this item.\par\medskip
\ifdefined\hypertarget\hypertarget{sourcepointer-5}{}\fi
\noindent\textit{Source reference.} Original reference 075S: sources/stacks-chow/tags.txt, source lines 9141--9141. Cross-chapter standard prerequisite cited by the author inside the counted span; it is used as a pointer and its own proof is not part of this item.\par\medskip
\begin{lemma}
\label{lemma-general-pushforward}
Let $\mathcal{M}$ be a rule which associates to every algebraic stack
$\mathcal{X}$ a subcategory $\mathcal{M}_\mathcal{X}$ of
$\textit{Mod}(\mathcal{X}_\etale, \mathcal{O}_\mathcal{X})$
such that
\begin{enumerate}
\item $\mathcal{M}_\mathcal{X}$ is a weak Serre subcategory
of $\textit{Mod}(\mathcal{X}_\etale, \mathcal{O}_\mathcal{X})$
(see Homology, Definition \ref{homology-definition-serre-subcategory})
for all algebraic stacks $\mathcal{X}$,
\item for a smooth morphism of algebraic stacks
$f : \mathcal{Y} \to \mathcal{X}$ the functor $f^*$ maps
$\mathcal{M}_\mathcal{X}$ into $\mathcal{M}_\mathcal{Y}$,
\item if $f_i : \mathcal{X}_i \to \mathcal{X}$ is a family of smooth
morphisms of algebraic stacks with
$|\mathcal{X}| = \bigcup |f_i|(|\mathcal{X}_i|)$, then an object
$\mathcal{F}$ of
$\textit{Mod}(\mathcal{X}_\etale, \mathcal{O}_\mathcal{X})$
is in $\mathcal{M}_\mathcal{X}$ if and only if
$f_i^*\mathcal{F}$ is in $\mathcal{M}_{\mathcal{X}_i}$ for all $i$, and
\item if $f : \mathcal{Y} \to \mathcal{X}$ is a morphism of algebraic
stacks such that $\mathcal{X}$ and $\mathcal{Y}$ are representable
by affine schemes, then $R^if_*$ maps $\mathcal{M}_\mathcal{Y}$
into $\mathcal{M}_\mathcal{X}$.
\end{enumerate}
Then for any quasi-compact and quasi-separated morphism 
$f : \mathcal{Y} \to \mathcal{X}$ of algebraic stacks
$R^if_*$ maps $\mathcal{M}_\mathcal{Y}$
into $\mathcal{M}_\mathcal{X}$. (Higher direct images computed in \'etale
topology.)
\end{lemma}

\begin{proof}
Let $f : \mathcal{Y} \to \mathcal{X}$ be a quasi-compact and quasi-separated
morphism of algebraic stacks and let $\mathcal{F}$ be an object of
$\mathcal{M}_\mathcal{Y}$. Choose a surjective smooth morphism
$\mathcal{U} \to \mathcal{X}$ where $\mathcal{U}$ is representable by
a scheme. By
Sheaves on Stacks, Lemma
\ref{stacks-sheaves-lemma-base-change-higher-direct-images}
taking higher direct images commutes with base change.
Assumption (2) shows that the pullback of $\mathcal{F}$
to $\mathcal{U} \times_\mathcal{X} \mathcal{Y}$ is in
$\mathcal{M}_{\mathcal{U} \times_\mathcal{X} \mathcal{Y}}$
because the projection
$\mathcal{U} \times_\mathcal{X} \mathcal{Y} \to \mathcal{Y}$
is smooth as a base change of a smooth morphism. Hence (3) shows we may
replace $\mathcal{Y} \to \mathcal{X}$ by the projection
$\mathcal{U} \times_\mathcal{X} \mathcal{Y} \to \mathcal{U}$.
In other words, we may assume that $\mathcal{X}$
is representable by a scheme.
Using (3) once more, we see that the question is Zariski local on
$\mathcal{X}$, hence we may assume that $\mathcal{X}$ is representable by
an affine scheme. Since $f$ is quasi-compact this implies that also
$\mathcal{Y}$ is quasi-compact. Thus we may choose a surjective smooth
morphism $g : \mathcal{V} \to \mathcal{Y}$ where $\mathcal{V}$ is representable
by an affine scheme.

\medskip\noindent
In this situation we have the spectral sequence
$$
E_2^{p, q} = R^q(f \circ g_p)_*g_p^*\mathcal{F}
\Rightarrow
R^{p + q}f_*\mathcal{F}
$$
of
Sheaves on Stacks, Proposition
\ref{stacks-sheaves-proposition-smooth-covering-compute-direct-image}.
Recall that this is a first quadrant spectral sequence hence we may
use the last part of Homology, Lemma \ref{homology-lemma-first-quadrant-ss}.
Note that the morphisms
$$
g_p : \mathcal{V}_p =
\mathcal{V} \times_\mathcal{Y} \ldots \times_\mathcal{Y} \mathcal{V}
\longrightarrow
\mathcal{Y}
$$
are smooth as compositions of base changes of the smooth morphism $g$.
Thus the sheaves $g_p^*\mathcal{F}$ are in
$\mathcal{M}_{\mathcal{V}_p}$ by (2). Hence it suffices to prove that the
higher direct images of objects of $\mathcal{M}_{\mathcal{V}_p}$ under
the morphisms
$$
\mathcal{V}_p =
\mathcal{V} \times_\mathcal{Y} \ldots \times_\mathcal{Y} \mathcal{V}
\longrightarrow
\mathcal{X}
$$
are in $\mathcal{M}_\mathcal{X}$. The algebraic stacks $\mathcal{V}_p$
are quasi-compact and quasi-separated by
Morphisms of Stacks, Lemma
\ref{stacks-morphisms-lemma-quasi-compact-quasi-separated-permanence}.
Of course each $\mathcal{V}_p$ is representable by an algebraic space
(the diagonal of the algebraic stack $\mathcal{Y}$ is representable
by algebraic spaces). This reduces us to the case where
$\mathcal{Y}$ is representable by an algebraic space and $\mathcal{X}$
is representable by an affine scheme.

\medskip\noindent
In the situation where $\mathcal{Y}$ is representable by an algebraic
space and $\mathcal{X}$ is representable by an affine scheme, we choose
anew a surjective smooth morphism $\mathcal{V} \to \mathcal{Y}$ where
$\mathcal{V}$ is representable by an affine scheme. Going through the
argument above once again we once again reduce to the morphisms
$\mathcal{V}_p \to \mathcal{X}$. But in the current situation the algebraic
stacks $\mathcal{V}_p$ are representable by quasi-compact and quasi-separated
schemes (because the diagonal of an algebraic space is representable by
schemes).

\medskip\noindent
Thus we may assume $\mathcal{Y}$ is representable by a scheme and
$\mathcal{X}$ is representable by an affine scheme. Choose (again)
a surjective smooth morphism $\mathcal{V} \to \mathcal{Y}$ where
$\mathcal{V}$ is representable by an affine scheme. In this case all
the algebraic stacks $\mathcal{V}_p$ are representable by separated
schemes (because the diagonal of a scheme is separated).

\medskip\noindent
Thus we may assume $\mathcal{Y}$ is representable by a separated scheme and
$\mathcal{X}$ is representable by an affine scheme. Choose (yet again)
a surjective smooth morphism $\mathcal{V} \to \mathcal{Y}$ where
$\mathcal{V}$ is representable by an affine scheme. In this case all
the algebraic stacks $\mathcal{V}_p$ are representable by affine schemes
(because the diagonal of a separated scheme is a closed immersion hence affine)
and this case is handled by assumption (4).
This finishes the proof.
\end{proof}

\end{document}
